\documentclass[11pt]{article}
\usepackage{amsmath,amsthm,amscd,amssymb,mathrsfs,tikz}
\usepackage{latexsym,epsf,epsfig}
\newcommand{\ds}{\displaystyle}
\newcommand{\sU}{\mathscr U}
\usepackage[left=1.2in, right=1.2in, top=1in, bottom=1in]{geometry}
%The above set the parameters adjust the margins. There are many ways to do this, and frankly, what is above is not the best. However, it will work the time being.
\usepackage{sectsty}


\usepackage{booktabs, multirow} % for borders and merged ranges
\usepackage{soul}% for underlines
\usepackage{xcolor,colortbl} % for cell colors
\usepackage{changepage,threeparttable} % for wide tables

\usepackage{hyperref}
\theoremstyle{plain}
\newtheorem{theorem}{Theorem}

\usepackage{graphicx}

% Margins
\topmargin=-0.45in
\evensidemargin=0in
\oddsidemargin=0in
\textwidth=6.5in
\textheight=9.0in
\headsep=0.25in

\title{ MATH 301: Homework 1}
\author{ Aren Vista }
\date{\today}

\begin{document}
\maketitle	
\pagebreak

\section*{Question 1: From Blackboard}
\textbf{Proposition:} Let $S$ and $T$ be two sets. If both $S$ and $T$ are countable, then $S \cup T$ is countable.
\textbf{Note:} If $S$ and $T$ are said to be infinite and countable; then they are said to be \textit{denumerable} 
\begin{proof} Suppose $S$ and $T$ are countable, and there exists $A,B \subset \mathbb{N}$, by definition there exists a bijective function:
    \[
        f: A \rightarrow S \text{ and } g: B \rightarrow T
    \] 
To show $S \cup T$ is countable, we must construct a bijection $h: C \rightarrow S \cup T, C \subset \mathbb{N}$ \\
\textbf{Define}  $h: C \rightarrow S \cup T$:
\[
    h(n) = 
    \begin{cases}
        f\left(\frac{n}{2}\right) & \text{if } n \text{ is even} \\
        g\left(\frac{n+1}{2}\right) & \text{if } n \text{ is odd} 
    \end{cases}
\] 
Note $S,T$ must be disjoint; else $\exists ~n_1, n_2 \in C, n_1 \neq n_2 \ni h(n_1) = f(n_1) = h(n_2) = g(n_2)$ \\
Thus, for the sake of the proof, presume $S,T$ to be disjoint sets s.t. $S \cap T = \emptyset$

\begin{center}
    To prove the proposition we must show that $h$ is \textit{surjective} and \textit{injective}. 
\end{center}

\subsection{Surjectivity}
\begin{proof} \textbf{Goal:} We want to show that $\forall ~y \in S \cup T, \exists ~n \in C \ni h(n) = y$
\subsubsection*{Case 1 ($y \in S$):}
Let $y \in S$. Since $f$ is surjective, $\exists ~ k \in C$ s.t. $f(k) = y$. \\
Let $n=2k$. By definition $n$ is even. Observe:
\[
    h(n) = f(\frac{2k}{2}) = f(k) = y
\] 
\subsubsection*{Case 2 ($y \in T$):}
Let $y \in T$. Since $g$ is surjective, $\exists ~k \in C \ni g(k) = y$ \\
Let $n=2k-1$. Then $n$ is odd. Observe:
\[
    h(n) = g(\frac{(2k-1)+1}{2}) = g(\frac{2k}{2}) = g(k) = y
\] 
As every element $y \in S \cup T \exists ~n \in C \ni h(n) = y \implies h: C \rightarrow S \cup T$ is surjective.  
\end{proof}
\pagebreak
\subsection{Injectivity}
\begin{proof} \textbf{Goal:} We want to show that for arbitrary elments $n_1,n_2 \in C$, $h(n1)=h(n_2) \implies n_1 = n_2$. Suppose $h(n_1) = h(n_2) = y$. 
\subsubsection*{Case 1 (Same Parity):}
Suppose $n_1, n_2$ be both even. Then the following statments hold.
\begin{align*}
    h(n_1) = h(n_2) & \implies f(\frac{n_1}{2}) = f(\frac{n_2}{2}). \\
    \text{Since } f \text{ is injective: } & \frac{n_1}{2} = \frac{n_2}{2} \implies n_1 = n_2
\end{align*}
Suppose $n_1, n_2$ be both odd. Then the following statments hold.
\begin{align*}
    h(n_1) = h(n_2) & \implies g(\frac{n_1 + 1}{2}) = g(\frac{n_2 + 1}{2}) \\
    \text{Since } g \text{ is injective: } & \frac{n_1 + 1}{2} = \frac{n_2 + 1}{2} \implies n_1 = n_2
\end{align*}
\subsubsection*{Case 2 (Mixed Parity):}
WLOG suppose $n_1$ is even and $n_2$ is odd. 
\begin{align*}
    & h(n_1) \in Img(f) = S \\
    & h(n_2) \in Img(g) = T \\
    \text{Since } S \cap T = \emptyset & \implies  h(n_1) \neq h(n_2) \equiv f(n_1) \neq g(n_2)
\end{align*}
As $n_1,n_2 \in C$, $h(n1)=h(n_2) \implies n_1 = n_2$. The function $h(n): C \rightarrow S \cup T$ is injective
\end{proof}
\subsection{Conclusion}
\begin{center}
Since $h$ is both surjective and injective, it is bijective. \\
As there exists a bijective function $h: C \rightarrow S \cup T \implies  S \cup T$ is countable. 
\end{center}
\end{proof}
\section*{Question 2: Sec 1.3 (q1)}
\textbf{Question:} Prove that a nonempty set $T_1$ is finite iff there is a bijection from $T_1$ onto a finite set $T_2$. 
\section*{Question 3: Sec 1.3 (q4)}
\textbf{Question:} Exhibit a bijection between $\mathbb{N}$ and the set of all odd integers greater than 13.
\section*{Question 4: Sec 1.3 (q6)}
\textbf{Question:} Exhibit a bijection between $\mathbb{N}$ and a proper subset of itself.
\section*{Question 5: Sec 1.3 (q7)}
\textbf{Question:} Prove that a set $T_1$ is denumerable iff there is a bijection from $T_1$ onto a denumerable set $T_2$
\section*{Question 6: Sec 1.3 (q8)}
\textbf{Question:} Give an example of a countable collection of finite sets whose union is not finite. \\
Consider the sets defined to be:
\[
    A_n = \{ n \}, \forall ~n \in \mathbb{N}
\] 
These sets are finite sets defined as:
\begin{align*}
    & A_1 = \{ 1 \} \\
    & A_2 = \{ 2 \} \\
    & A_3 = \{ 3 \} \\
    & ... \\
    & A_n = \{ n \}
\end{align*}
Observe each set $A_1, A_2, A_3, ..., A_2$ are finite (singleton) sets. \\
\begin{center}
Though the union of all sets, by how we defined $A_1, A_2, A_3, ..., A_2$, would be result in the equivalent set $\mathbb{N}$ which is not finite.
\end{center}

\end{document}
